\documentclass[11pt, a4paper]{article}
\usepackage[margin=2.5cm]{geometry}
\usepackage{booktabs}
\usepackage{array}
\usepackage{longtable}
\usepackage{xcolor}
\usepackage{hyperref}
\usepackage{microtype}
\usepackage{parskip}
\usepackage{enumitem}
\usepackage{graphicx}
\usepackage{caption}
\usepackage{subcaption}
\usepackage[T1]{fontenc}
\usepackage[utf8]{inputenc}
\usepackage{tabularx}
\usepackage{multirow}
\usepackage{hhline}
\usepackage{amsmath}
\usepackage{adjustbox}

% ── Global table defaults ────────────────────────────────────────────────────
% Tighten inter-column spacing and add a touch of row breathing room
\setlength{\tabcolsep}{4pt}
\renewcommand{\arraystretch}{1.08}
% New column type W: paragraph cell that fills remaining width in tabularx
\newcolumntype{W}{>{\raggedright\arraybackslash}X}

\definecolor{highprio}{HTML}{CC0000}
\definecolor{medprio}{HTML}{CC8800}
\definecolor{lowprio}{HTML}{007700}
\definecolor{good}{HTML}{1A9850}
\definecolor{mid}{HTML}{CC8800}
\definecolor{poor}{HTML}{CC0000}
\definecolor{rowgray}{HTML}{F5F5F5}

\hypersetup{
  colorlinks = true,
  linkcolor  = blue!60!black,
  urlcolor   = blue!60!black,
}

\newcommand{\good}[1]{{\color{good}\textbf{#1}}}
\newcommand{\poor}[1]{{\color{poor}#1}}
\newcommand{\na}{---}

\title{\textbf{ERBOT v0.2.0 --- Integrated Benchmark Results}\\[6pt]
       \large CORA \quad$\cdot$\quad Affiliation Strings \quad$\cdot$\quad
              Synthetic 10K \quad$\cdot$\quad NCVR 5-Party \quad$\cdot$\quad NCVR 10-Party}
\author{ERBOT v0.2.0}
\date{Generated: 2026-03-26}

\begin{document}
\maketitle
\tableofcontents
\vspace{1em}
\hrule
\vspace{1em}

% ═════════════════════════════════════════════════════════════════════════════
\section*{Abstract}
\addcontentsline{toc}{section}{Abstract}

This document integrates the individual ERBOT benchmark results across five
datasets: CORA (bibliographic deduplication), Affiliation Strings (free-text
institution deduplication), Synthetic 10K (noisy unstructured text
deduplication), NCVR 5-Party (two-source voter record linkage at two scales),
and NCVR 10-Party (multi-source voter record linkage at 100K scale). It
synthesises per-dataset findings into cross-cutting conclusions about method
reliability, blocking effectiveness, scalability limits, and recommended
configurations for ERBOT v0.2.0.

% ═════════════════════════════════════════════════════════════════════════════
\section{Benchmark Overview}

\subsection{Datasets}

Six experimental conditions across five benchmark datasets were evaluated.

\begin{table}[h!]
\centering
\small
\begin{adjustbox}{max width=\linewidth}
\begin{tabular}{llrrrrl}
\toprule
\textbf{ID} & \textbf{Dataset} & \textbf{Records} & \textbf{Fields} &
\textbf{Gold pairs} & \textbf{Match rate$^*$} & \textbf{Task} \\
\midrule
\textbf{C}    & CORA Bibliographic       & 1,879   & 15 & 64,578  & 3.7\%   & Deduplication \\
\textbf{A}    & Affiliation Strings      & 2,260   &  1 & 32,816  & 1.3\%   & Deduplication \\
\textbf{S}    & Synthetic 10K (d10k)     & 10,000  &  1 & 8,705   & 0.017\% & Deduplication \\
\textbf{N5-s} & NCVR 5-Party (small)     & 10,000  &  4 & $\sim$880 & 0.018\% & Linkage \\
\textbf{N5-l} & NCVR 5-Party (large)     & 100,000 &  4 & $\sim$880 & $<$0.001\% & Linkage \\
\textbf{N10}  & NCVR 10-Party            & 100,000 &  4 & 1,512   & 0.0003\% & Linkage \\
\bottomrule
\end{tabular}
\end{adjustbox}
\caption{Benchmark conditions. $^*$Match rate = gold pairs / $\binom{n}{2}$ (fraction of all possible
record pairs that are true matches). CORA and Affiliation are relatively dense;
single-field and linkage datasets are very sparse.}
\end{table}

\subsection{Pipeline Configuration Summary}

\begin{table}[h!]
\centering
\small
\begin{tabular}{llrrrl}
\toprule
\textbf{ID} & \textbf{Blocking} & \textbf{Pairs} & \textbf{RR} &
\textbf{Methods run} & \textbf{Runtime} \\
\midrule
\textbf{C}    & Prefix-3 on \texttt{title}       & 137,707   & 92.2\% & 9 (incl.\ SVM)      & $\approx$59 min \\
\textbf{A}    & Prefix-3 on \texttt{affiliation} & 164,719   & 93.5\% & 9 (incl.\ SVM)      & $\approx$5.6 hr \\
\textbf{S}    & Prefix-3 on \texttt{text}        & 885,301   & 98.2\% & 8 (SVM excluded$^a$)  & $\approx$11 min \\
\textbf{N5-s} & Prefix-3 on \texttt{surname}     & $\sim$10,600 & $\approx$100\% & 9 (incl.\ SVM) & $<$1 hr \\
\textbf{N5-l} & Prefix-3 on \texttt{surname}     & 1,059,972$^b$ & $\approx$100\% & 4 (graph-only$^c$) & $\approx$45 hr \\
\textbf{N10}  & Prefix-3 on \texttt{surname}     & 5,275,664 & 99.9\% & 3 (graph-only$^c$) & $\approx$8 min \\
\bottomrule
\end{tabular}
\caption{%
$^a$ SVM killed after $\approx$20 hours (885K pairs, estimated 1--4 days to complete).
$^b$ Truncated from 2,121,188 (pair budget cap exceeded).
$^c$ hclust/pam/gc excluded: exceed R's 65,536-observation distance-matrix limit or
require 74.5\,GiB dense allocation. N10 uses 5M pair budget vs.\ 2M for N5-l.}
\end{table}

% ═════════════════════════════════════════════════════════════════════════════
\section{Per-Dataset Results}

\subsection{CORA Bibliographic Dataset (C)}

\textbf{Characteristics:} 1,879 research paper citations with 15 fields; heavy
noise (OCR errors, missing fields, formatting inconsistencies).
Ground truth: 64,578 matched pairs.

\begin{table}[h!]
\centering
\small
\setlength{\tabcolsep}{4pt}
\begin{tabular}{lrrrrrr}
\toprule
\textbf{Method} & \textbf{ARI} & \textbf{NMI} & \textbf{B$^3$-P} & \textbf{B$^3$-R} & \textbf{B$^3$-F} & \textbf{PF-F} \\
\midrule
\good{hclust\_avg} & \good{0.746} & \good{0.855} & \good{0.651} & 0.849 & \good{0.737} & \good{0.757} \\
hclust\_ward     & 0.524 & 0.827 & 0.661 & 0.748 & 0.702 & 0.542 \\
pam              & 0.624 & 0.841 & 0.639 & 0.784 & 0.704 & 0.638 \\
svm              & 0.543 & 0.849 & 0.630 & 0.817 & 0.711 & 0.566 \\
threshold\_cc    & 0.506 & 0.811 & 0.520 & 0.858 & 0.648 & 0.532 \\
louvain          & 0.506 & 0.811 & 0.520 & 0.858 & 0.648 & 0.532 \\
label\_prop      & 0.506 & 0.811 & 0.520 & 0.858 & 0.648 & 0.532 \\
leiden           & 0.000 & 0.686 & 1.000 & 0.065 & 0.122 & 0.000 \\
gc$^\dagger$     & 0.000 & 0.000 & 0.040 & 1.000 & 0.077 & 0.076 \\
\midrule
\textit{final}   & 0.746 & 0.855 & 0.651 & 0.849 & 0.737 & 0.757 \\
\bottomrule
\end{tabular}
\caption{CORA performance (rerun 2026-03-26, merge bug fixed).
$^\dagger$gc failed with dimension mismatch error.}
\end{table}

\textbf{Key points:} hclust\_avg (ARI = 0.746) is the best method on CORA, and the
highest ARI result across all deduplication benchmarks. PAM (0.624) and SVM (0.543) follow.
The \texttt{final} merge now correctly selects hclust\_avg. With 15 discriminative fields,
hierarchical and centroidal methods outperform the graph-based family (threshold-CC: 0.506).

\subsection{Affiliation Strings Dataset (A)}

\textbf{Characteristics:} 2,260 author affiliation strings, one noisy text
field, extreme surface-form variability (abbreviations, language variants,
reorderings). Ground truth: 32,816 matched pairs.

\begin{table}[h!]
\centering
\small
\setlength{\tabcolsep}{4pt}
\begin{tabular}{lrrrrrr}
\toprule
\textbf{Method} & \textbf{ARI} & \textbf{NMI} & \textbf{B$^3$-P} & \textbf{B$^3$-R} & \textbf{B$^3$-F} & \textbf{PF-F} \\
\midrule
\good{threshold\_cc} & \good{0.076} & 0.688 & 0.395 & 0.550 & \good{0.460} & \na \\
louvain              & 0.076 & 0.688 & 0.394 & 0.550 & 0.459 & \na \\
label\_prop          & 0.076 & 0.688 & 0.394 & 0.550 & 0.459 & \na \\
svm                  & 0.035 & 0.818 & 0.977 & 0.159 & 0.273 & \na \\
leiden               & 0.000 & \good{0.819} & 1.000 & 0.146 & 0.255 & \na \\
hclust\_avg          & 0.053 & 0.516 & 0.184 & 0.530 & 0.272 & \na \\
pam                  & 0.053 & 0.516 & 0.184 & 0.530 & 0.272 & \na \\
hclust\_ward         & 0.042 & 0.516 & 0.185 & 0.519 & 0.273 & \na \\
gc                   & 0.000 & 0.000 & 0.007 & 1.000 & 0.014 & \na \\
\midrule
\textit{final}       & 0.076 & 0.688 & 0.395 & 0.550 & 0.460 & \na \\
\bottomrule
\end{tabular}
\caption{Affiliation Strings performance. PF columns omitted (in individual summary).}
\end{table}

\textbf{Key points:} Overall performance is very low (best ARI = 0.076).
SVM paradoxically underperforms threshold-CC (ARI 0.035 vs.\ 0.076): extreme
precision (0.977) but near-zero recall (0.159) --- it only merges the highest-confidence
pairs. Primary bottleneck: prefix-3 blocking on a variable-opening-token field
discards many true matches before any clustering is attempted. Final merge
correctly selected threshold-CC.

\subsection{Synthetic 10K Dataset (S)}

\textbf{Characteristics:} 10,000 noisy unstructured text records, one field,
corruptions via character substitution, token deletion, and reordering.
Ground truth: 8,705 matched pairs.

\begin{table}[h!]
\centering
\small
\setlength{\tabcolsep}{4pt}
\begin{tabular}{lrrrrrr}
\toprule
\textbf{Method} & \textbf{ARI} & \textbf{NMI} & \textbf{B$^3$-P} & \textbf{B$^3$-R} & \textbf{B$^3$-F} & \textbf{PF-F} \\
\midrule
leiden               & 0.000 & \good{0.914} & \good{1.000} & 0.271 & \good{0.426} & \na \\
\good{threshold\_cc} & \good{0.023} & 0.778 & 0.334 & 0.519 & 0.406 & \na \\
louvain              & 0.023 & 0.778 & 0.334 & 0.519 & 0.406 & \na \\
label\_prop          & 0.023 & 0.778 & 0.334 & 0.519 & 0.406 & \na \\
hclust\_avg          & 0.002 & 0.371 & 0.019 & 0.573 & 0.036 & \na \\
pam                  & 0.002 & 0.371 & 0.019 & 0.573 & 0.036 & \na \\
hclust\_ward         & 0.001 & 0.347 & 0.019 & 0.690 & 0.037 & \na \\
gc                   & 0.000 & 0.000 & 0.001 & 1.000 & 0.002 & \na \\
svm                  & \multicolumn{6}{c}{\textit{not run --- estimated 1--4 days; killed after 20 hr}} \\
\midrule
\textit{final}       & 0.023 & 0.778 & 0.334 & 0.519 & 0.406 & \na \\
\bottomrule
\end{tabular}
\caption{Synthetic 10K performance. PF columns omitted (in individual summary).}
\end{table}

\textbf{Key points:} Performance is the lowest of any single-condition run (best ARI = 0.023).
An interesting inversion: Leiden has the highest NMI (0.914) and B$^3$-F (0.426) but ARI = 0
due to over-splitting. Threshold-CC, Louvain, and Label Propagation tie at ARI = 0.023 with more
balanced recall (0.519 vs.\ 0.271). SVM is infeasible. Blocking (98.2\% reduction) is again
the primary bottleneck.

\subsection{NCVR 5-Party Dataset (N5)}

\textbf{Characteristics:} North Carolina Voter Registration records, 4 fields
(\texttt{givenname}, \texttt{surname}, \texttt{suburb}, \texttt{postcode}).
The dataset file contains 5 parties; the current runs used parties 0 and 1
(two-source linkage), 20\% overlap, 2 fields modified per duplicate.
Run at two scales: 5K/party (10K total, small) and 50K/party (100K total, large).

\begin{table}[h!]
\centering
\small
\begin{adjustbox}{max width=\linewidth}
\begin{tabular}{llrrrrrr}
\toprule
\textbf{Scale} & \textbf{Method} & \textbf{ARI} & \textbf{NMI} & \textbf{B$^3$-P} &
\textbf{B$^3$-R} & \textbf{B$^3$-F} & \textbf{PF-F} \\
\midrule
\multirow{5}{*}{\textbf{Small}}
 & \good{threshold\_cc / louvain /} & \good{0.917} & \good{0.975} & \good{1.000} & 0.929 & \good{0.963} & \good{0.923} \\
 & \good{\quad label\_prop / svm (tied)} & & & & & & \\
 & hclust\_avg & 0.276 & 0.713 & 0.537 & 0.929 & 0.681 & 0.364 \\
 & pam         & 0.158 & 0.672 & 0.537 & 0.786 & 0.638 & 0.258 \\
 & leiden      & 0.000 & 0.849 & 1.000 & 0.500 & 0.667 & 0.000 \\
 & gc          & 0.000 & 0.000 & 0.143 & 1.000 & 0.250 & 0.143 \\
\midrule
\multirow{5}{*}{\textbf{Large}}
 & \good{threshold\_cc} & \good{0.195} & \good{0.956} & \good{0.991} & 0.560 & \good{0.715} & \good{0.195} \\
 & louvain     & 0.133 & 0.947 & 0.933 & 0.560 & 0.700 & 0.134 \\
 & label\_prop & 0.133 & 0.947 & 0.933 & 0.560 & 0.700 & 0.134 \\
 & leiden      & 0.000 & 0.951 & 1.000 & 0.500 & 0.667 & 0.000 \\
 & hclust/pam/gc & \multicolumn{6}{c}{\textit{failed (65K limit / 74.5 GiB allocation)}} \\
\midrule
\multirow{1}{*}{}
 & \textit{final (both)} & 0.917 / 0.195 & 0.975 / 0.956 & 1.000 / 0.991 & 0.929 / 0.560 & 0.963 / 0.715 & 0.923 / 0.195 \\
\bottomrule
\end{tabular}
\end{adjustbox}
\caption{NCVR 5-Party performance at small (10K records) and large (100K records) scale.}
\end{table}

\textbf{Key points:} NCVR at small scale is the strongest result across all
benchmarks (ARI = 0.917), benefiting from a 4-field structured profile and
controlled synthetic noise. At large scale, ARI drops to 0.195 --- not because
the clustering degrades, but because the match density is 10$\times$ lower
(880 true matches in 100K records vs.\ 10K records). Precision stays at 0.991;
recall drops from 0.929 to 0.560. Four methods (hclust\_avg, hclust\_ward, pam,
gc) fail at 100K scale and have been removed from the large-run script.

% ─────────────────────────────────────────────────────────────────────────────
\subsection{NCVR 10-Party Dataset (N10)}

\textbf{Characteristics:} 100,000 NCVR records across 10 parties (10K per party);
4 structured fields; 20\% overlap, 2-field modification. True cross-party pairs:
1,512 (match density 0.030\%). Graph-only methods; 5M pair budget.

\begin{table}[h!]
\centering
\small
\begin{adjustbox}{max width=\linewidth}
\begin{tabular}{lrrrrrr}
\toprule
\textbf{Method} & \textbf{B$^3$-P} & \textbf{B$^3$-R} & \textbf{B$^3$-F} &
\textbf{Hom.} & \textbf{Com.} & \textbf{V-measure} \\
\midrule
\textbf{threshold\_cc} & \textbf{0.824} & \textbf{0.991} & \textbf{0.900} & \textbf{0.951} & \textbf{0.999} & \textbf{0.975} \\
louvain               & 0.629 & 0.991 & 0.769 & 0.835 & 0.999 & 0.909 \\
label\_prop           & 0.629 & 0.991 & 0.769 & 0.835 & 0.999 & 0.909 \\
\midrule
\textit{final}        & 0.824 & 0.991 & 0.900 & 0.951 & 0.999 & 0.975 \\
\bottomrule
\end{tabular}
\end{adjustbox}
\caption{NCVR 10-Party performance (100K records). ARI, NMI, VI, and PairF not computed
at this scale (O($n^2$) evaluation bottleneck). Hom.\ = Homogeneity; Com.\ = Completeness.}
\end{table}

\textbf{Key points:} threshold\_cc outperforms louvain/label\_prop by a large margin
(B$^3$-F 0.900 vs.\ 0.769) in contrast to NCVR5-small where all three methods tied.
At very low match density (1,512 true pairs in 100K records), threshold\_cc's strict
cut-off suppresses spurious merges more effectively than community-detection methods.
The 5M pair budget (vs.\ 2M in NCVR5-large) is a key difference: NCVR5-large budget
truncation likely contributed to its lower recall (0.560), while NCVR10 with 5M pairs
achieves near-perfect recall (0.991). Note that NCVR5-large and NCVR10 are not directly
comparable (different party counts and true pair distributions) --- this is directional
evidence, not a controlled budget experiment.
Total runtime was 8.2 minutes (vs.\ $\approx$45 hours for NCVR5-large) --- the
improvement is entirely due to removing failed distance-matrix methods.

% ═════════════════════════════════════════════════════════════════════════════
\section{Cross-Dataset Analysis}

\subsection{Overall Performance Comparison}

\begin{table}[h!]
\centering
\small
\begin{adjustbox}{max width=\linewidth}
\begin{tabular}{lrrrrrrl}
\toprule
\textbf{Dataset} & \textbf{Best ARI} & \textbf{Best B$^3$-F} & \textbf{Best V} &
\textbf{Best PF-F} & \textbf{Final B$^3$-F} & \textbf{Final correct?} & \textbf{Best method} \\
\midrule
CORA          & 0.746 & 0.737 & 0.855 & 0.757 & 0.737 & Yes & hclust\_avg \\
Affiliation   & 0.076 & 0.460 & 0.688 & \na   & 0.460 & Yes & threshold\_cc \\
Syn10K        & 0.023 & 0.426 & 0.914 & \na   & 0.406 & Yes & threshold\_cc \\
NCVR5 (small) & 0.917 & 0.963 & 0.975 & 0.923 & 0.963 & Yes & threshold\_cc \\
NCVR5 (large) & 0.195 & 0.715 & 0.956 & 0.195 & 0.715 & Yes & threshold\_cc \\
NCVR10        & \na   & 0.900 & 0.975 & \na   & 0.900 & Yes & threshold\_cc \\
\bottomrule
\end{tabular}
\end{adjustbox}
\caption{Best performance per dataset/condition. ``Final correct'' = did the \texttt{best}
merge strategy select the actual best method? NCVR10 ARI/PairF not available (evaluation
bottleneck at $n=100{,}000$). V = V-measure.}
\end{table}

\subsection{Method Consistency Across Datasets}

\begin{table}[h!]
\centering
\small
\begin{tabularx}{\linewidth}{lrrrrrrW}
\toprule
\textbf{Method} & \textbf{CORA} & \textbf{Aff} & \textbf{Syn10K} &
\textbf{N5-s} & \textbf{N5-l} & \textbf{N10$^\dagger$} & \textbf{Assessment} \\
\midrule
\textbf{hclust\_avg}    & \good{0.746} & 0.053 & 0.002 & 0.276 & fail  & n/a  & Best on CORA; weak elsewhere; fails $>$65K \\
\textbf{threshold\_cc} & 0.506 & 0.076 & 0.023 & 0.917 & 0.195 & \good{0.900\,B$^3$} & Best or tied-best in all 6 conditions \\
\textbf{louvain}       & 0.506 & 0.076 & 0.023 & 0.917 & 0.133 & 0.769\,B$^3$ & Matches threshold-CC except NCVR-large \\
\textbf{label\_prop}   & 0.506 & 0.076 & 0.023 & 0.917 & 0.133 & 0.769\,B$^3$ & Identical to louvain in all conditions \\
pam                    & 0.624 & 0.053 & 0.002 & 0.158 & fail  & n/a  & Strong on CORA; weak elsewhere; fails $>$65K \\
svm                    & 0.543 & 0.035 & n/a   & 0.917 & n/a   & n/a  & Good on multi-field; infeasible at scale \\
hclust\_ward           & 0.524 & 0.042 & 0.001 & 0.068 & fail  & n/a  & Moderate on CORA; weak elsewhere; fails $>$65K \\
leiden                 & 0.000 & 0.000 & 0.000 & 0.000 & 0.000 & n/a  & ARI = 0 in all conditions \\
gc                     & 0.000 & 0.000 & 0.000 & 0.000 & 0.000 & n/a  & ARI = 0 / fails in all conditions \\
\bottomrule
\end{tabularx}
\caption{Primary metric by method and condition (n/a = not run; fail = runtime error).
CORA/Aff/Syn10K/N5 values are ARI; N10$^\dagger$ values are B$^3$-F (ARI unavailable at 100K scale).}
\end{table}

Three clear groups emerge:

\begin{enumerate}[leftmargin=*, label=\arabic*., noitemsep]
  \item \textbf{Reliable graph methods:} threshold\_cc, louvain, label\_prop are best
        or joint-best in every condition. They converge to an identical solution in 4 of
        5 ARI-comparable conditions. In the N10 (B$^3$-F) run, threshold\_cc pulls ahead
        of louvain/label\_prop by 0.131 B$^3$-F at very low match density.
        threshold\_cc is the single most robust method and should be the ERBOT default.
  \item \textbf{Conditionally useful:} SVM wins on CORA (ARI +0.037 over graph methods) where
        15 fields provide rich training signal. On single-field or large datasets it is worse
        than threshold-CC or infeasible. The SVM viability boundary appears to be
        $\approx$200K candidate pairs (complete in reasonable time up to 164K pairs, fails at
        885K).
  \item \textbf{Never useful at defaults:} leiden (ARI = 0 across all 5 ARI-comparable
        conditions), gc (ARI = 0 across all 5 ARI-comparable conditions), hclust\_avg/ward,
        pam. These methods either degenerate or fail at scale without parameter tuning.
        They should not appear in production runs without explicit tuning.
\end{enumerate}

\subsection{Blocking Effectiveness}

\begin{table}[h!]
\centering
\small
\begin{tabularx}{\linewidth}{lrlrlW}
\toprule
\textbf{Dataset} & \textbf{RR} & \textbf{Blocking key} & \textbf{Gold pairs} &
\textbf{Bottleneck?} & \textbf{Notes} \\
\midrule
CORA         & 92.2\% & Prefix-3 on \texttt{title}       & 64,578  & Moderate & Title prefix usually consistent in citations \\
Affiliation  & 93.5\% & Prefix-3 on \texttt{affiliation} & 32,816  & \textbf{Severe}   & Variable opening tokens (``Dept.\ of'', ``Univ.\ of'') \\
Syn10K       & 98.2\% & Prefix-3 on \texttt{text}        & 8,705   & \textbf{Severe}   & Corrupted/reordered first tokens \\
NCVR5-small  & $\approx$100\% & Prefix-3 on \texttt{surname}   & $\sim$880 & Low & Surname prefix stable under 2-field modification \\
NCVR5-large  & $\approx$100\% & Prefix-3 on \texttt{surname}   & $\sim$880 & \textbf{High}$^*$ & Budget cap (2.12M $\to$ 2M) truncated true pairs; recall=0.560 \\
NCVR10       & 99.9\% & Prefix-3 on \texttt{surname}     & 1,512   & Low & 5M budget sufficient; recall=0.991; no truncation \\
\bottomrule
\end{tabularx}
\caption{Blocking effectiveness per dataset. $^*$NCVR5-large 2M budget was the primary
cause of recall=0.560; raising to 5M (as in NCVR10) recovers recall=0.991.}
\end{table}

Prefix-3 blocking on the primary text field is appropriate only when the field
has a consistent, reliable prefix (as in bibliographic titles and personal surnames).
It is a poor choice for free-text affiliation strings and synthetic noisy text,
where many true match pairs have incompatible opening tokens. For structured records
(NCVR), surname prefix is stable and the limiting factor is the pair budget, not
prefix quality.

\subsection{Scalability Profile}

\begin{table}[h!]
\centering
\small
\begin{tabular}{lrll}
\toprule
\textbf{Method} & \textbf{Max viable records} & \textbf{Bottleneck} & \textbf{Alternative} \\
\midrule
hclust\_avg/ward & $<$65,536 & $O(n^2)$ distance matrix & graph methods \\
pam              & $<$65,536 & $O(n^2)$ distance matrix & graph methods \\
gc               & $<$65,536 & Dense matrix allocation   & remove from defaults \\
svm              & $\sim$200K pairs & Super-linear kernel computation & logistic regression, GBM \\
threshold\_cc    & $\geq$100,000 & Sparse graph ops & --- \\
louvain          & $\geq$100,000 & Sparse graph ops & --- \\
label\_prop      & $\geq$100,000 & Sparse graph ops & --- \\
leiden           & $\geq$100,000 & Sparse graph ops & --- \\
\bottomrule
\end{tabular}
\caption{Scalability limits by method.}
\end{table}

\subsection{Performance Drivers}

The following factors explain performance variation across conditions,
in approximate order of impact:

\begin{enumerate}[leftmargin=*, label=\arabic*.]
  \item \textbf{Field count and noise quality} (primary). NCVR with 4 structured fields
        and controlled noise achieves ARI = 0.917; the single-field datasets plateau at
        0.023--0.076. CORA with 15 fields but heavy uncontrolled noise (OCR errors,
        missing fields) achieves ARI = 0.746 --- lower than NCVR despite having more
        fields. Field count helps when fields are informative; noise quality determines
        how discriminative each field is. Both factors drive performance jointly.

  \item \textbf{Match density.} At the same 4-field NCVR configuration, match density
        decreases as scale grows: small (N5-s, $\sim$880 pairs in 10K records) $\to$
        large (N5-l, $\sim$880 in 100K) drops ARI from 0.917 to 0.195. When true
        matches are rare, blocking recall failures and false-positive cluster merges
        are amplified.

  \item \textbf{Blocking recall.} On single-field noisy text (Affiliation, Syn10K),
        prefix-3 blocking is too aggressive: many true match pairs have incompatible
        prefixes and are eliminated before any similarity or clustering computation.
        The measured performance is a lower bound for these datasets.

  \item \textbf{Supervision signal.} SVM outperforms graph methods on CORA (+0.037 ARI)
        where 15 fields provide a rich multi-dimensional feature space. On single-field
        datasets the feature space is one-dimensional and SVM's decision boundary
        collapses to a scalar threshold --- functionally equivalent to threshold-CC but
        slower.

  \item \textbf{Noise model.} Controlled 2-field corruption (NCVR) is easier than
        uncontrolled real-world errors (CORA citations) or extreme surface-form variation
        (Affiliation). Datasets with predictable noise are easier for all methods.
\end{enumerate}

\subsection{The \texttt{final} Merge Strategy}

\begin{table}[h!]
\centering
\begin{tabular}{lllrr}
\toprule
\textbf{Dataset} & \textbf{Selected} & \textbf{Correct?} & \textbf{Final ARI} & \textbf{Final B$^3$-F} \\
\midrule
CORA          & hclust\_avg   & \good{Yes} & 0.746 & 0.737 \\
Affiliation   & threshold\_cc & \good{Yes} & 0.076 & 0.460 \\
Syn10K        & threshold\_cc & \good{Yes} & 0.023 & 0.406 \\
NCVR5-small   & threshold\_cc & \good{Yes} & 0.917 & 0.963 \\
NCVR5-large   & threshold\_cc & \good{Yes} & 0.195 & 0.715 \\
NCVR10        & threshold\_cc & \good{Yes} & \na   & 0.900 \\
\bottomrule
\end{tabular}
\caption{Final merge selection per condition. \na = not computed (evaluation bottleneck).
The merge bug was present in the original CORA run and fixed before subsequent runs.}
\end{table}

The merge selection bug (originally exposed by CORA) has been fixed. All six
conditions now correctly identify the best method. The CORA rerun (2026-03-26)
confirms \texttt{final} = hclust\_avg at ARI = 0.746 --- the strongest ARI result
across all benchmarks. NCVR10 achieves B$^3$-F = 0.900 with threshold\_cc, the
highest B-cubed score of any large-scale (100K) condition.

% ═════════════════════════════════════════════════════════════════════════════
\section{Unified Key Findings}

\begin{enumerate}[leftmargin=*, label=\arabic*.]

  \item \textbf{Threshold-CC is the most reliable ERBOT method.}
        It is best or tied-best in all six experimental conditions, scales to
        $\geq$100K records, and is fast. It should be the primary default
        clustering method for all ERBOT runs. At very low match density
        (NCVR10, 0.030\%), threshold\_cc separates from louvain/label\_prop
        by 0.131 B$^3$-F points --- its strict cut-off criterion matters more
        as true match pairs become rare.

  \item \textbf{Louvain and Label Propagation are robust second choices.}
        They produce results identical to threshold-CC in 4 of 5 ARI-comparable conditions.
        At very large scale with low match density (NCVR5-large, NCVR10) they
        underperform threshold-CC. All three graph methods are worth running
        together since their agreement is a strong quality signal.

  \item \textbf{Leiden and GC should be excluded from default runs.}
        Both produce ARI = 0 in every condition where ARI was available without parameter
        tuning. Leiden over-partitions (too-high resolution leading to near-zero recall)
        and GC over-merges (degenerate fallback caused by an unresolved dimension mismatch
        bug). Until tuned per-dataset, they add noise to the method comparison and inflate
        runtime.

  \item \textbf{hclust\_avg is the best method on CORA (ARI = 0.746), but does not
        generalise to other datasets.}
        On CORA's dense 15-field similarity matrix, average-link hierarchical clustering
        finds an excellent partition. On single-field datasets or at $>$65K records it fails
        or degenerates. SVM (ARI = 0.543 on CORA) is more portable: it also ties for best
        on NCVR-small, though it is infeasible above $\approx$200K candidate pairs.

  \item \textbf{The \texttt{final} merge is working correctly across all datasets.}
        The original CORA run had a merge selection bug that picked the worst method.
        After fixing and rerunning (2026-03-26), CORA \texttt{final} correctly selects
        hclust\_avg at ARI = 0.746. All six conditions now report correct \texttt{final}
        results.

  \item \textbf{Blocking is the primary performance bottleneck on single-field and
        noisy datasets.}
        Prefix-3 blocking is well-suited to bibliographic titles and personal surnames,
        but discards large fractions of true matches on affiliation strings and
        synthetic noisy text. Pair completeness (fraction of gold pairs retained after
        blocking) is not currently measured and should be added to Stage 3 diagnostics.

  \item \textbf{Distance-matrix methods (hclust, pam, gc) are not viable at scale.}
        They fail hard at $>$65K records due to R's memory constraints. They should
        be gated behind a record-count check in \texttt{er\_cluster\_all()} and
        automatically excluded when $n > 65000$.

  \item \textbf{Performance correlates strongly with field count.}
        NCVR (4 fields, controlled noise) $\gg$ CORA (15 fields, uncontrolled noise)
        $\gg$ Affiliation/Syn10K (1 field). Every additional discriminative field
        substantially reduces the difficulty of the clustering task.

  \item \textbf{Pair budget is the dominant NCVR performance lever at scale.}
        NCVR5-large with a 2M budget achieves recall 0.560; NCVR10 with a 5M budget
        achieves recall 0.991 at the same $n=100{,}000$ scale. Raising the budget from
        2M to 5M recovers 0.431 recall points. Budget should scale with record count.

  \item \textbf{Equal weights are reasonable but not optimal.}
        CORA has 15 fields of highly varying quality; supervised weight learning would
        likely add +0.05--0.10 ARI for SVM and threshold-CC. For NCVR (4 balanced
        fields) and single-field datasets, weight learning provides no benefit.

  \item \textbf{Evaluation scalability is a bottleneck at $n \geq 100{,}000$.}
        ARI, NMI, VI, and PairF could not be computed for NCVR10 due to O($n^2$)
        loops in \texttt{er\_bcubed()} and contingency-table routines. B-cubed and
        V-measure completed. This should be fixed with vectorised table lookups before
        any further large-scale benchmarking.

\end{enumerate}

% ═════════════════════════════════════════════════════════════════════════════
\section{Unified Recommendations}

\begin{table}[h!]
\centering
\small
\begin{tabular}{cp{0.80\linewidth}}
\toprule
\textbf{Priority} & \textbf{Action} \\
\midrule
{\color{highprio}\textbf{High}} &
  \textbf{Fix O($n^2$) evaluation bottleneck} in \texttt{er\_bcubed()} and the
  contingency-table routines. Replace inner \texttt{which()} loops with vectorised
  table lookups so ARI, NMI, VI, and PairF can be computed at $n = 100{,}000$. \\[4pt]
{\color{highprio}\textbf{High}} &
  \textbf{Gate distance-matrix methods behind a record-count check.}
  In \texttt{er\_cluster\_all()}, automatically skip hclust\_avg, hclust\_ward, pam, and gc
  when $n > 65000$ rather than attempting and failing. \\[4pt]
{\color{highprio}\textbf{High}} &
  \textbf{Add pair completeness to Stage 3 output.}
  Report the fraction of ground-truth pairs retained after blocking. Currently all
  recall figures are lower bounds. \\[4pt]
{\color{highprio}\textbf{High}} &
  \textbf{Remove leiden and gc from default method lists},
  or add a degenerate-detection guard that flags ARI = 0 results and excludes
  them from the \texttt{final} merge comparison. \\[4pt]
{\color{medprio}\textbf{Medium}} &
  \textbf{Replace SVM with logistic regression or gradient boosting}
  for datasets with $>$200K candidate pairs. Provides near-linear scalability
  with similar accuracy in the multi-field setting. \\[4pt]
{\color{medprio}\textbf{Medium}} &
  \textbf{Improve blocking for single-field and free-text datasets.}
  For affiliation strings: sorted neighbourhood (SN) blocking or token-set
  Jaccard. For synthetic noisy text: SN blocking on the full field or n-gram
  shingling. \\[4pt]
{\color{medprio}\textbf{Medium}} &
  \textbf{Set pair budget to at least 5$\times$N} for structured linkage tasks.
  NCVR10 demonstrated that a 5M budget at 100K records achieves recall=0.991,
  while the old 2M budget caused recall=0.560. Make the budget adaptive or
  scale with record count. \\[4pt]
{\color{medprio}\textbf{Medium}} &
  \textbf{Tune Leiden resolution parameter per dataset.}
  Default resolution is consistently too high; try 0.05--0.2. Until tuned,
  exclude Leiden from production comparisons. \\[4pt]
{\color{medprio}\textbf{Medium}} &
  \textbf{Fix GC dimension mismatch bug.}
  GC silently falls back to a degenerate solution on all datasets where it was run
  due to an unresolved interface error. Either fix or remove from the method list. \\[4pt]
{\color{lowprio}\textbf{Low}} &
  \textbf{Implement supervised weight learning for multi-field datasets.}
  ARI-based weight optimisation on CORA expected to gain +0.05--0.10 ARI over
  equal weights. \\[4pt]
{\color{lowprio}\textbf{Low}} &
  \textbf{Use all 5 parties in the NCVR5 dataset} (currently only parties 0 and 1
  are used) to assess 5-source linkage under the full NCVR5 design.
  The 10-party design (NCVR10) has already been run; this fills the gap at
  intermediate party count. \\
\bottomrule
\end{tabular}
\end{table}

% ═════════════════════════════════════════════════════════════════════════════
\section{Conclusions}

ERBOT v0.2.0 has been evaluated across five benchmark datasets (six conditions)
spanning a wide range of ER task types and sizes. The results establish a clear
picture of what works, what scales, and what needs fixing.

\textbf{What works reliably:} Threshold-CC (and by extension Louvain and Label
Propagation) is a robust, scalable default that performs best or jointly best
in every condition. At very low match density (NCVR10, 0.030\%), threshold\_cc
pulls 0.131 B$^3$-F ahead of community-detection methods, making it even more
valuable at large scale. The \texttt{final} merge strategy correctly identifies
the best method in all six conditions (the CORA merge bug is now fixed).

\textbf{What the numbers mean in context:} The best achievable ARI varies
enormously by task --- from 0.917 on NCVR5-small (controlled noise, 4 fields)
and 0.746 on CORA (15 fields, rich signal) to 0.023 on Syn10K (noisy single
field). NCVR10 (B$^3$-F = 0.900) is the strongest large-scale result and shows
that structured 4-field linkage is tractable at 100K records when the pair
budget is sized appropriately. These differences reflect fundamentally different
task difficulties, not ERBOT's limitations.

\textbf{What needs immediate attention:} Three issues stand out. First,
evaluation scalability: ARI, NMI, VI, and PairF could not be computed for
NCVR10 due to O($n^2$) inner loops --- fixing this is a prerequisite for
further large-scale benchmarking. Second, the four degenerate methods (leiden,
gc, hclust/pam at scale) generate runtime failures and misleading zero-ARI
entries in every result table --- they should be gated or removed. Third, pair
completeness at Stage 3 is unmeasured, meaning all recall figures are lower
bounds.

\textbf{Direction:} The most impactful improvements would come from fixing
evaluation scalability (to fully characterise NCVR10 and future large-scale
runs), better blocking for free-text fields (Affiliation, Syn10K), and a
scalable replacement for SVM. The NCVR10 result confirms that graph-only ERBOT
runs at 100K scale are fast ($<$10 minutes) and accurate (B$^3$-F $\approx$0.90)
when the pair budget is adequate.

% ═════════════════════════════════════════════════════════════════════════════
\section{Output Files Reference}

\begin{table}[h!]
\centering
\small
\begin{tabular}{ll}
\toprule
\textbf{Path} & \textbf{Contents} \\
\midrule
\texttt{res\_cora/CORA\_RESULTS\_SUMMARY.pdf}     & Full CORA individual summary \\
\texttt{res\_cora/er\_performance.csv}            & CORA metric table \\
\texttt{res\_aff/AFF\_RESULTS\_SUMMARY.pdf}       & Full Affiliation individual summary \\
\texttt{res\_aff/er\_performance.csv}             & Affiliation metric table \\
\texttt{res\_syn10k/SYN10K\_RESULTS\_SUMMARY.pdf} & Full Syn10K individual summary \\
\texttt{res\_syn10k/er\_performance.csv}          & Syn10K metric table \\
\texttt{res\_ncvr5/NCVR5\_RESULTS\_SUMMARY.pdf}   & Full NCVR5 individual summary \\
\texttt{res\_ncvr5/small/er\_performance.csv}     & NCVR5-small metric table \\
\texttt{res\_ncvr5/large/er\_performance.csv}     & NCVR5-large metric table \\
\texttt{res\_ncvr10/NCVR10\_RESULTS\_SUMMARY.pdf} & Full NCVR10 individual summary \\
\texttt{res\_ncvr10/10k/er\_performance.csv}      & NCVR10 metric table \\
\texttt{ERBOT\_INTEGRATED\_SUMMARY.pdf}           & This document \\
\bottomrule
\end{tabular}
\end{table}

\end{document}
